\firstsection{Motivation}
\maketitle

\vspace{-1mm}

Uncertainty visualization has been recognized as one of the grand research challenges that must be addressed for trusted scientific discovery and decision-making~\cite{TA:2004:Johnson:topScivis,TA:Johnson:2003:nextStepVisErrors}. Visualization of uncertainty, however, is nontrivial given many difficult relevant challenges, including uncertainty modeling, computation, communication, perception, cognition, and decision-making under uncertainty. Furthermore, widespread use of artificial intelligence (AI) technologies across scientific workflows makes visualization and analysis of uncertainty even more challenging. Although the VIS conference brings together experts from subdomains related to uncertainty, including visualization and high-performance computing, large part of the community does not directly work on uncertainty visualization. Furthermore, it is paramount to create a venue that can attract experts outside the visualization field, including application science, applied math, AI, and psychology, to be able to tackle theoretical and application-specific uncertainty visualization challenges. Our workshop, therefore, will provide a platform that would foster interdisciplinary interaction and dissemination of work, which will be critical to identifying research priorities for tackling uncertainty visualization challenges that are currently obstacles to scientific progress. 

%The unique aspect that sets our workshop apart from the main VIS conference is its potential to bring together experts outside and within various branches of the visualization domain to solve challenges in uncertainty. Arranging such workshop is absolutely essential to help push research in uncertainty visualization in right direction. 




% Uncertainty visualization has been recognized as one of the grand research challenges that must be addressed for trusted scientific discovery and decision-making~\cite{TA:2004:Johnson:topScivis,TA:Johnson:2003:nextStepVisErrors}. Visualization of uncertainty is however nontrivial given many difficult challenges, including computation, rendering, perception, and cognition, and decision-making under uncertainty. Although the visualization community has experts in separate sub-domains, including computation, rendering, perception, and cognition, the large part of the community has not directly worked on uncertainty visualization. Bringing these experts together is therefore crucial to addressing the tough challenges in uncertainty analysis. Further, creating a venue that can attract contributions from experts outside the visualization field (e.g., application science, applied math, human perception and cognition) is also important to enhancing theoretical and application-specific understanding of uncertainty. Our workshop therefore aims to set up a platform that would promote interdisciplinary interaction and dissemination of work to advance uncertainty visualization research. The unique aspect that sets our workshop apart from the main VIS conference is its potential to bring together experts outside and within various branches of the visualization domain to solve challenges in uncertainty. Arranging such workshop is absolutely essential to help push research in uncertainty visualization in right direction.    

%Addressing these challenges will need expertise and contributions from diverse backgrounds not necessarily confined to the visualization domain. For example, understanding the importance and treatment of uncertainty in application domains would require engagement with application scientists who may not have a background in visualization. Further, the interaction with experts in mathematics, statistics, high-performance computing, and cognitive science is crucial to understanding and progressing computation, perception, and decision-making aspects of uncertainty. Our workshop therefore aims to set up a platform that would promote interdisciplinary 
%interaction and dissemination of work to advance uncertainty visualization research. The unique aspect that sets our workshop apart from the main VIS conference is its potential to attract exciting submissions from domain experts, mathematicians, and psychologists who may not be the core visualization researchers. On the other hand, arranging such workshop is absolutely essential to the visualization community to help push research in uncertainty visualization in right direction.    
%\textbf{Dave: Instead of specifically naming psychologist, how about: domain experts in application science, computational science, applied math, human perception and cognition? I modified the abstract to fit these ideas.}

 %\ta{invite domain-experts who have never been to vis before}
 %\textbf{Dave: Question on the title- Is 'scientific' too limiting? Would eliminating the word cast a wider net? SciVis is a small part of the Vis conference.  Would be interesting to try and bring in more of the wider vis community. Could be some opportunities for cross-fertilization of ideas and collaboration.}
%\ta{Explain why a separate uncertainty workshop is needed? Why and what things could be missed by the main VIS (and LDAV, TopoInVis etx) event that will be addressed by this workshop}.